\documentclass{article}
\usepackage[T1]{fontenc}
\usepackage{lmodern}
\usepackage{graphicx}
\usepackage{amsmath,amssymb}
\usepackage[a4paper,margin=2cm]{geometry}
\usepackage[absolute,overlay]{textpos}
\setlength{\TPHorizModule}{1mm}
\setlength{\TPVertModule}{1mm}
\usepackage[indonesian]{babel}
\usepackage{hyperref}
\setlength{\emergencystretch}{2em}
% unit: Halaman 4

\begin{document}

% === Halaman 4 (llm) ===

\section*{Properti Algoritma RSA}

\begin{enumerate}
    \item $p$ dan $q$ bilangan prima \hfill \textcolor{red}{(rahasia)}
    \item $n = p \cdot q$ \hfill (tidak rahasia)
    \item $\phi(n) = (p - 1)(q - 1)$ \hfill \textcolor{red}{(rahasia)}
    \item $e$ (kunci enkripsi) \hfill (tidak rahasia, disebut kunci publik)

    Syarat: $\text{PBB}(e, \phi(n)) = 1$, PBB = pembagi bersama terbesar = $gcd$

    \item $d$ (kunci dekripsi) \hfill \textcolor{red}{(rahasia, disebut kunci privat)}

    $d$ dihitung dari $d \equiv e^{-1} \pmod{\phi(n)}$

    \item $m$ (plainteks) \hfill \textcolor{red}{(rahasia)}
    \item $c$ (cipherteks) \hfill (tidak rahasia)
\end{enumerate}

\end{document}
